% math4cs-hw11-group.tex

% !TEX program = xelatex
%%%%%%%%%%%%%%%%%%%%
% see http://mirrors.concertpass.com/tex-archive/macros/latex/contrib/tufte-latex/sample-handout.pdf
% for how to use tufte-handout
\documentclass[a4paper, justified]{tufte-handout}

\input{hw-preamble} % feel free to modify this file if you understand LaTeX well
%%%%%%%%%%%%%%%%%%%%
\title{计算机数学 (12-groups: 群论)}
\me{魏恒峰}{hfwei@hnu.edu.cn}{}{}
\date{2026年6月11日}
%%%%%%%%%%%%%%%%%%%%
\begin{document}
\maketitle
%%%%%%%%%%%%%%%%%%%%
\noplagiarism % PLEASE DON'T DELETE THIS LINE!
%%%%%%%%%%%%%%%%%%%%
\begin{abstract}
\end{abstract}
%%%%%%%%%%%%%%%%%%%%
\beginrequired
%%%%%%%%%%%%%%%

%%%%%%%%%%%%%%%
\begin{problem}
  在整数集 $\mathbb{Z}$ 中, 规定运算 $\oplus$ 如下:
  \[
    \forall a, b \in \mathbb{Z}\; a \oplus b = a + b - 2.
  \]
  请证明: $(\mathbb{Z}, \oplus)$ 构成群。
\end{problem}

\begin{proof}
\end{proof}
%%%%%%%%%%%%%%%

%%%%%%%%%%%%%%%
\begin{problem}
  设 $G$ 是群。
  请证明: 如果 $\forall x \in G\; x^2 = e$,
  则 $G$ 是交换群。
\end{problem}

\begin{proof}
\end{proof}
%%%%%%%%%%%%%%%

%%%%%%%%%%%%%%%
\begin{problem}
  请求出 $3^{83}$ 的最后两位数~\footnote{\url{https://www.wolframalpha.com/input/?i=3\%5E83}}。要求给出计算过程。
\end{problem}

\begin{proof}
\end{proof}
%%%%%%%%%%%%%%%

%%%%%%%%%%%%%%%
\begin{problem}
  设 $H$ 是群 $G$ 的子集，即 $H \le G$。
  请证明: 对任意的 $g \in G$, 集合
  \[
    gHg^{-1} = \{ghg^{-1} \mid h \in H\}
  \]
  是 $G$ 的子群。
\end{problem}

\begin{proof}
\end{proof}
%%%%%%%%%%%%%%%

%%%%%%%%%%%%%%%
\begin{problem}
  设 $\phi$ 是群 $G$ 到群 $G'$ 的同构映射。
  请证明:
  \[
    \forall g \in G\; \text{ord } g = \text{ord } \phi(g)\text{。}
  \]
\end{problem}

\begin{proof}
\end{proof}
%%%%%%%%%%%%%%%

%%%%%%%%%%%%%%%
\begin{problem}
  请给出循环群 $(\mathbb{Z}_{18}, +)$ 的所有生成元与所有子群。
\end{problem}

\begin{solution}
\end{solution}
%%%%%%%%%%%%%%%

%%%%%%%%%%%%%%%
\begin{problem}
  设正整数 $n$ 的质因数分解为
  \[
    n = p_{1}^{\alpha_{1}} p_{2}^{\alpha_{2}} \cdots p_{k}^{\alpha_{k}}\text{。}
  \]
  其中, $p_{1}, p_{2}, \cdots, p_{k}$ 是 $n$ 的所有不同的质因数。\\[3pt]

  \noindent 请证明: 循环群 $(\mathbb{Z}_{n}, +)$ 的子群的个数为
  \[
    (\alpha_{1} + 1)(\alpha_{2} + 1) \cdots (\alpha_{k} + 1)\text{。}
  \]
\end{problem}

\begin{proof}
\end{proof}
%%%%%%%%%%%%%%%

%%%%%%%%%%%%%%%%%%%%
% 如果没有需要订正的题目，可以把这部分删掉
% \begincorrection
%%%%%%%%%%%%%%%%%%%%

%%%%%%%%%%%%%%%%%%%%
% 如果没有反馈，可以把这部分删掉
\beginfb

你可以写 (也可以发邮件)
\begin{itemize}
  \item 对课程及教师的建议与意见
  \item 教材中不理解的内容
  \item 希望深入了解的内容
  \item $\cdots$
\end{itemize}
%%%%%%%%%%%%%%%%%%%%
\end{document}